\begin{table}[t]
\caption{Исследовательские вопросы и данные. Полный экспериментальный протокол сохранён в дополнительных материалах.}\label{tab:questions}
\centering\small\setlength{\tabcolsep}{3pt}\begin{tabular}{p{.25\textwidth}p{.32\textwidth}p{.34\textwidth}}\toprule
Вопрос & Данные & Проверяемое утверждение\\\midrule
Выбор модели поведения & Зафиксированный опыт подготовки до кризиса & Статическая модель, отбор по прогнозу и сохранение нагрузочных компонент выбирают разные действия.\\
Источник различий & Анализ правил, простой выбор и планирование по истинному генератору & Разделение выбора действия, порога и ошибки модели.\\
Условия применения & Чувствительность к эффекту, критерию события и горизонту; проверка выявления структуры & Установление условий, охваченных оценкой реакции.\\
Наблюдаемый сигнал нагрузки & Прогноз на открытых журналах BPI & Нагрузка может помогать прогнозу; реальные вмешательства здесь не оцениваются.\\\bottomrule
\end{tabular}\end{table}

\subsection{Диагностический анализ и вспомогательные данные}
После завершения основного эксперимента мы зафиксировали отдельный диагностический протокол. Он добавляет простое правило: сравнить оценённое усиление длительности и повторной работы между 10-м и 90-м процентилями нагрузки в обучающих данных, затем выбрать соответствующее действие при пороге 20. Оценки равны логарифмам отношений ожидаемой длительности B и ожидаемого множителя числа повторений $1/(1-p)$; отрицательные значения заменяются нулём, сложность фиксируется равной нулю. При большей оценке повторной работы выбирается её сокращение, иначе --- добавление ресурса. Правило использует только обученные параметры и состояния из обучающих данных.

Планировщик с известным генератором применяет исходные сценарии и повторы планирования к истинным уравнениям. Его выбор сохраняется до оценки по существующим независимым тестовым действиям. Это эталон знания модели в конечном каталоге правил, а не глобально оптимальный контроллер с доступом к будущему. Отдельно определяется лучшее фиксированное правило на реализованных тестовых данных, явно обозначенное ретроспективной нижней границей. Разности стоимости с этими правилами измеряют разные источники потерь решения. Интервалы последующего анализа описательные и не являются новыми подтверждающими критериями.

Проверка эффекта вмешательства сохраняет исходные решения, принятые при предполагаемом сокращении возвратов на 50\%, и оценивает фактические сокращения на 30\% и 70\%. Другие проверки используют сохранённые траектории при порогах незавершённых задач 80/100/120, длительностях события 100/200/300, горизонтах 1400/1600/1800 и весах конечного остатка задач 0/100/200. Меняются условия оценки, но правила не выбираются по тестовым исходам. Проверка выявления структуры использует только обучающие трассы и явно отделяет восстановленное последовательное семейство от прежних ошибок выявления параллельных процессов.

Прежние симуляционные кампании оценивают прогноз траекторий и структурную точность. Открытые журналы BPI 2012 и BPI 2017~\cite{bpi12,bpi17} дают хронологические проверки признаков нагрузки для прогноза длительности и следующего старта. Они выявляют наблюдаемый сигнал, а не подтверждают эффекты вмешательств; полные методы, контрольные сравнения и результаты приведены в дополнительных материалах.
